\section*{अध्याय \ref{ch:b1:work-and-energy} --- कार्य, ऊर्जा और स्थितिज ऊर्जा}

\begin{solution}{exo:b1:work-and-energy:1}
उठाने वाला बल $+mgh = 20 \times 9.81 \times 3.0 = \qty{589}{J}$; भार
$-\qty{589}{J}$; शक्ति $589/5.0 = \qty{118}{W}$। घर्षण: $-\mu_dmgd =
-0.30 \times 196 \times 3.0 = \qty{-177}{J}$।
\end{solution}

\begin{solution}{exo:b1:work-and-energy:2}
$E_k = \tfrac12 \times 1000 \times 27.8^2 = \qty{386}{kJ}$; $d = E_k/F =
\qty{55}{m}$; $v$ दुगुना करने पर $E_k$ और $d$ चौगुने: \qty{220}{m}।
\end{solution}

\begin{solution}{exo:b1:work-and-energy:3}
$E_p = \tfrac12 \times 500 \times 0.10^2 = \qty{2.5}{J}$; $v = \sqrt{2 \times
2.5/0.050} = \qty{10}{m/s}$।
\end{solution}

\begin{solution}{exo:b1:work-and-energy:4}
$\tfrac12 mv^2 = mg\ell(1 - \cos 60^\circ)$: $v = \sqrt{g\ell} = \qty{3.1}{m/s}$।
तनाव: $T - mg = mv^2/\ell = mg$, अतः $T = 2mg$।
\end{solution}

\begin{solution}{exo:b1:work-and-energy:5}
घर्षण $-\mu_dmg\cos\alpha\,(h/\sin\alpha)$ \omterm{def:b1:work-and-energy:work}{कार्य} करता है, अतः
\[
\tfrac12 mv^2 = mgh(1 - \mu_d\cot\alpha) = mgh(1 - 0.35),
\]
$v = \sqrt{2 \times 9.81 \times 2.0 \times 0.65} = \qty{5.1}{m/s}$; $35\%$ खो गया।
\end{solution}

\begin{solution}{exo:b1:work-and-energy:6}
पिंड के अनंत तक पहुँचने के लिए $E_m = \tfrac12 mv^2 - GMm/R \geq 0$ चाहिए
(जहाँ $E_p = 0$ और $v \to 0$): $v_e = \sqrt{2GM/R}$, अर्थात्
\[
v_e = \sqrt{\frac{2 \times 6.67\times10^{-11} \times 5.97\times10^{24}}{6.37\times10^6}}
= \qty{11.2}{km/s}.
\]
समान घनत्व पर $M \propto R^3$: अतः $v_e \propto R$, यानी \qty{22.4}{km/s}।
\end{solution}

\begin{solution}{exo:b1:work-and-energy:7}
निम्नतम बिंदु: $mg(\ell_0 + x) = \tfrac12 kx^2$, $25x^2 - 687x - 13734 = 0$,
$x = \qty{41}{m}$; कुल पतन \qty{61}{m}; $T_{\max} = kx = \qty{2.0}{kN}
= 3.0\,mg$ (शुद्ध \qty{2}{g} ऊपर की ओर)। चाल वहाँ सबसे अधिक है जहाँ शुद्ध
बल लुप्त हो जाता है, $kx = mg$: $x = \qty{14}{m}$, यानी छलाँग से \qty{34}{m} नीचे।
\end{solution}

\begin{solution}{exo:b1:work-and-energy:8}
गुरुत्व: $mgv\sin\alpha = 80 \times 9.81 \times 5.0 \times 0.05 = \qty{196}{W}$;
कर्षण: $\tfrac12\rho C_xSv^3 = 0.5 \times 1.2 \times 0.50 \times 125 =
\qty{38}{W}$; कुल \qty{234}{W}।
\end{solution}

\begin{solution}{exo:b1:work-and-energy:9}
$E_p' = 4E_0(x^3/a^4 - x/a^2) = 0$: $x = 0, \pm a$। $E_p'' = E_0(12x^2/a^4 -
4/a^2)$: $-4E_0/a^2$, जो $0$ पर है (अस्थायी उच्चिष्ठ, $E_p = 0$), और $+8E_0/a^2$,
जो $\pm a$ पर है (स्थायी निम्निष्ठ, $E_p = -E_0$): यानी $E_0$ गहराई के कूप;
$\omega_0^2 = 8E_0/ma^2$।
\end{solution}

\begin{solution}{exo:b1:work-and-energy:10}
$\tfrac12 m\dot x^2 + \tfrac12 kx^2 = E$: यह $x_{\max}
= \sqrt{2E/k}$ और $\dot x_{\max} = \sqrt{2E/m}$ अर्ध-अक्षों वाला दीर्घवृत्त
है। गति आवर्त है और उसका आवर्तकाल $2\pi\sqrt{m/k}$ है, \omterm{def:b1:wave-propagation:signal}{आयाम} चाहे जो हो
(समकालिकता), अतः हर दीर्घवृत्त उतना ही समय लेता है।
\end{solution}

\begin{solution}{exo:b1:work-and-energy:11}
ऊर्जा: $v^2 = 2gR(1 - \cos\theta)$। त्रिज्य दिशा: $mg\cos\theta - N = mv^2/R$,
अतः $N = mg(3\cos\theta - 2)$, जो $\cos\theta = 2/3$ पर शून्य है ($\theta =
48^\circ$): ऊँचाई $2R/3$, चाल $\sqrt{2gR/3}$।
\end{solution}

\begin{solution}{exo:b1:work-and-energy:12}
$\tfrac12 mv^2 = \mu_dmgd$: $\mu_d = 625/(2 \times 9.81 \times 80) = 0.40$।
ऊर्जा टायर--सड़क के संपर्क में ऊष्मा बन गई (और एक घसीट-निशान)। लुढ़कते
पहिये स्थैतिक घर्षण काम में लेते हैं, $\mu_s > \mu_d$: बड़ा ब्रेक-बल, कम
दूरी, और दिशा-नियंत्रण भी बना रहता है।
\end{solution}

\begin{solution}{pb:b1:work-and-energy:1}
\textbf{1.} $r \to 0$: $+\infty$ ($r^{-12}$ पद के कारण); $r \to \infty$:
$0^-$। एक खड़ी दीवार, एक कूप, और एक लंबी पूँछ।

\textbf{2.} $E_p' = \epsilon[-12\sigma^{12}/r^{13} + 12\sigma^6/r^7] = 0$ तभी जब
$(\sigma/r)^6 = 1$: $r = \sigma$, $E_p = \epsilon(1 - 2) = -\epsilon$।

\textbf{3.} $(\sigma/r)^{12} = u^2$: $E_p = \epsilon(u^2 - 2u) = \epsilon[(u - 1)^2
- 1]$, जो $u = 1$ ($r = \sigma$) पर न्यूनतम है और उसका मान $-\epsilon$ है।

\textbf{4.} $F = -E_p' = (12\epsilon/\sigma)[(\sigma/r)^{13} - (\sigma/r)^7]$:
$r < \sigma$ के लिए धनात्मक (प्रतिकर्षी), और उससे आगे ऋणात्मक (आकर्षी)।

\textbf{5.} $(1/1.2)^6 = 0.335$, $(1/1.2)^{12} = 0.112$: $E_p = \epsilon(0.112
- 0.670) = -0.56\epsilon = \qty{-2.5}{eV}$। $(1/1.2)^{13} = 0.0935$,
$(1/1.2)^7 = 0.279$: $F = (12\epsilon/\sigma)(-0.186) = -2.2\epsilon/\sigma =
-2.2 \times 7.0\times10^{-19}/1.27\times10^{-10} = \qty{-12}{nN}$।

\textbf{6.} $\epsilon = \qty{4.4}{eV} = 4.4 \times 96.5 = \qty{425}{kJ/mol}$,
जो मापे गए \qty{431}{kJ/mol} के निकट है।

\textbf{7.} $E_p'' = \epsilon[156\sigma^{12}/r^{14} - 84\sigma^6/r^8]$; और
$\sigma$ पर: $72\epsilon/\sigma^2 > 0$: अतः स्थायी।

\textbf{8.} $E_p \approx -\epsilon + \tfrac12(72\epsilon/\sigma^2)(r - \sigma)^2$:
अतः $k = 72\epsilon/\sigma^2$।

\textbf{9.} $k = 72 \times 7.0\times10^{-19}/(1.27\times10^{-10})^2 =
\qty{3.1e3}{N/m}$।

\textbf{10.} $\omega_0 = \sqrt{k/m} = \sqrt{3.1\times10^3/1.67\times10^{-27}} =
\qty{1.4e15}{rad/s}$; $f = \qty{2.2e14}{Hz}$; $\lambda = c/f = \qty{1.4}{\micro m}$;
$1/\lambda = \qty{7300}{cm^{-1}}$।

\textbf{11.} $2.4$ गुना अधिक: प्रतिरूप की दीवारें बहुत खड़ी हैं
($r^{-12}$ क्रोड सहसंयोजक बंध का एक भोंडा विकल्प है), पर \omterm{def:b1:units-dimensions:oom}{परिमाण की कोटि}
--- $10^{14}$ Hz पर एक अवरक्त कंपन --- ठीक है, और विधि (कूप की वक्रता)
ठीक वही है जो बेहतर स्थितिजों के साथ भी काम में आती है।

\textbf{12.} $\tfrac12 kA^2 \approx k_BT$: $A = \sqrt{2k_BT/k} = \sqrt{8.3\times
10^{-21}/3.1\times10^3} = \qty{1.6e-12}{m} = 0.013\sigma$: बंध बस काँपता
भर है।

\textbf{13.} $v_{\max} = A\omega_0 = 1.6\times10^{-12} \times 1.4\times10^{15}
= \qty{2.2}{km/s}$।

\textbf{14.} $-\epsilon < E_m < 0$: $(\sigma, 0)$ के चारों ओर बंद वक्र ---
कंपन, बद्ध; $E_m > 0$: खुले वक्र --- परमाणु अलग हो जाते हैं;
$E_m = 0$: विभाजक वक्र, यानी वह परमाणु जो बस अनंत तक पहुँच पाता है।

\textbf{15.} $u^2 - 2u = -\tfrac12$, जहाँ $u = (\sigma/r)^6$: $u = 1 \pm 1/\sqrt2
= 1.71$ या $0.293$; $r = \sigma u^{-1/6}$: यानी $0.915\sigma$ और $1.23\sigma$।

\textbf{16.} मध्यबिंदु $1.07\sigma > \sigma$: बाहरी दीवार भीतरी से कोमल है,
इसलिए परमाणु बाहर की ओर अधिक जगह पाता है; कालमाध्य दूरी ऊर्जा के साथ बढ़ती
है --- गरम बंध लंबे होते हैं, ठोस फैलते हैं।

\textbf{17.} बड़ा: कोमल बाहरी ओर प्रत्यानयन बल आवर्त सन्निकटन की मान्यता से
दुर्बल है, इसलिए लौटने में अधिक समय लगता है।

\textbf{18.} $\epsilon/2 = \qty{2.2}{eV} = \qty{3.5e-19}{J}$; $\sigma$ पर
\omterm{thm:b1:work-and-energy:ke}{गतिज ऊर्जा} के रूप में: $v = \sqrt{2 \times 3.5\times10^{-19}/1.67\times10^{-27}}
= \qty{2.0e4}{m/s}$।

\textbf{19.} ऊर्जा संरक्षण: $E_m > 0$ रहता है $> 0$, वक्र खुला रहता है,
और परमाणु एक ही भेंट के बाद उड़ जाते हैं। कोई तीसरा पिंड (दूसरा अणु, कोई
दीवार, कोई फ़ोटॉन) अतिरिक्त ऊर्जा ले जाना ही चाहिए।

\textbf{20.} $W = E_p(2\sigma) - E_p(\sigma) = \epsilon(2^{-12} - 2 \times 2^{-6})
+ \epsilon = 0.97\epsilon = \qty{4.3}{eV}$, धनात्मक: आकर्षण परमाणु को भीतर
खींचता है, अतः प्रेरक।

\textbf{21.} $W_{\text{बाह्य}} = E_p(1.5\sigma) - E_p(\sigma) = \epsilon(0.0077 -
0.176) + \epsilon = 0.83\epsilon = \qty{3.7}{eV}$।

\textbf{22.} $(1/0.9)^{13} = 3.93$, $(1/0.9)^7 = 2.09$: $F = +22\epsilon/\sigma
= \qty{+120}{nN}$, यानी $1.2\sigma$ पर के आकर्षण का दस गुना: किसी बंध को
$10\%$ दबाना उसे $20\%$ तानने से कहीं महँगा है।

\textbf{23.} $T = \epsilon/k_B = 7.0\times10^{-19}/1.38\times10^{-23} =
\qty{5.1e4}{K}$। असली गैसें इससे कहीं नीचे वियोजित हो जाती हैं, क्योंकि
ऊर्जाएँ वितरित होती हैं: अणुओं का एक अंश $k_BT$ से कई गुना ऊर्जा लिए रहता
है (\cref{ch:b1:kinetic-theory})।

\textbf{24.} बचे रहेंगे: निम्निष्ठ पर साम्यावस्था, निम्निष्ठ पर
$k = E_p''$, वियोजन ऊर्जा $=$ गहराई, $E_p = E_m$ से \omterm{def:b1:work-and-energy:turning}{पलटाव बिंदु}, और विषमता
$\Rightarrow$ प्रसार। बदलेंगे: $k$ के, आवृत्ति के और पलटाव बिंदुओं के मान।

\textbf{25.} साम्यावस्था: $E_p' = 0$; कंपन: $\omega_0^2 = E_p''/m$;
पलायन: $E_m$ की कूप की कोर से तुलना।
\end{solution}
